\section{Zusammenfassung}

%Durch dieses Projekt hatten wir die Möglichkeit, unser Wissen und unseren Umgang mit Software Engineering praktisch anzuwenden. Während der Vorlesungszeit haben wir uns mit wesentlichen Themen beschäftigt, die uns dabei geholfen haben, eine strukturierte Webanwendung zu entwickeln und die Anforderungen systematisch zu erfassen. Mithilfe von Personas, Anforderungsanalysen, Use-Cases sowie Aktivitätsdiagrammen konnten wir die wichtigsten Anforderungen der Anwendung darstellen und anschließend in die Implementierung überführen. Durch die Kombination von Servlets, MariaDB und Redis konnten wir die Geschäftslogik umsetzen, Datenbankzugriffe effizient gestalten und eine Warteschlange für die Druckaufträge realisieren. Zusätzlich haben wir ein Design Pattern ausgewählt, analysiert und in unserer Anwendung umgesetzt. Dabei entschieden wir uns für das \textbf{Singleton Pattern}, welches beispielsweise für die zentrale Verwaltung bestimmter Ressourcen eingesetzt werden kann. Durch die Verwendung von vorbereiteten SQL-Anweisungen konnten wir außerdem Sicherheitsrisiken wie \textbf{SQL-Injection} vermeiden. Darüber hinaus beschäftigten wir uns mit den von Balzert beschriebenen Analysemustern und deren UML-Darstellungen. Dabei konnten wir mögliche Probleme innerhalb eines Systems erkennen und passende Lösungsansätze anhand der Muster untersuchen. Um die Qualität unseres Quellcodes sicherzustellen, integrierten wir einen Git-Hook, der vor jedem Push überprüft, ob der Java-Code den festgelegten Formatierungsregeln entspricht. Dadurch konnten Fehler bereits frühzeitig erkannt und verhindert werden, dass fehlerhafter Code in das gemeinsame Repository gelangt. Ein weiterer wichtiger Bestandteil des Projekts war die Durchführung von Lasttests. Hierfür entwickelten wir Testszenarien, welche typische Benutzeraktionen wie Anmeldung, Ausleihe und Abmeldung simulieren. Die daraus gewonnenen Messwerte wurden anschließend visualisiert und analysiert, um das Verhalten der Anwendung unter paralleler Belastung besser beurteilen zu können. Durch die Umsetzung des Projekts konnten wir außerdem wichtige Erfahrungen im Bereich der \textit{Qualitätssicherung} sammeln. Wir haben gelernt, dass Softwarequalität nicht nur durch funktionierende Implementierungen entsteht, sondern durch ein Zusammenspiel verschiedener Maßnahmen wie automatisierten Tests, Code-Qualitätsprüfungen, Versionsverwaltung, Sicherheitsbetrachtungen und Performanceanalysen. Besonders die Integration von Prüfmechanismen in den Entwicklungsprozess hat uns gezeigt, wie Fehler frühzeitig erkannt und die Wartbarkeit einer Anwendung langfristig verbessert werden kann.
%Zusammenfassend konnten wir durch dieses Projekt viele theoretische Inhalte aus dem Bereich Software Engineering praktisch anwenden und ein besseres Verständnis dafür entwickeln, wie komplexe Softwaresysteme geplant, entwickelt, getestet und kontinuierlich verbessert werden.

Durch dieses Projekt konnten wir wesentliche Konzepte des Software Engineerings
praktisch anwenden. Mithilfe von Personas, Anforderungsanalysen, dem Use-Case und dem
Aktivitätsdiagram wurden die Anforderungen systematisch erfasst und in eine
strukturierte Webanwendung überführt. Die Kombination aus Servlets, MariaDB und
Redis ermöglichte eine effiziente Umsetzung der Geschäftslogik sowie eine
Druckwarteschlange nach dem Producer-Consumer-Muster. Zur Qualitätssicherung
integrierten wir einen Git-Hook, der vor jedem Push die Java-Formatierung prüft,
sowie Lasttests, die das Systemverhalten unter paralleler Belastung analysieren.
Das Projekt hat gezeigt, dass Softwarequalität durch das Zusammenspiel von
Implementierung, Sicherheit, Versionsverwaltung und Performanceanalyse entsteht.
